\scanpage{4}
% L1–5
en posant
\[
S_i=\Delta\times\{i\}\qquad,
\]
et on a un iso. canonique $\omega_{S_i}\simeq\omega_\Delta$ pour \unsure{tt|$\forall$} $i\in\unsure{\tau|T}$, d'où
\[
\bigwedge_{i\in\tau}\omega_{S_i}\simeq\unsure{\mathbb{1}|\underline{1}}.
\quad\text{(torseur unité sous $\mathbb Z/2\mathbb Z$)}
\]

% L6–12
Considérons la structure \deleted{donnée} sur un ens $S$ de card $6$ \unsure{donnée|donné} par une partition
\[
S=\coprod_{i\in\tau}S_i
\]
du type $(3,3)$, et un \inserted{\unsure{syst. trans.} d'}\unsure{isom|isom.} entre les $\omega_{S_i}$ i.e. un isom
\[
\bigwedge_{i\in\tau}\omega_{S_i}\simeq\mathbb{1}.\tag*{$(*)$}
\]

% L13–15
\deleted{Considérons les} Le groupe des \unsure{autom.|automorph.} de cette structure est d'ordre $36$ (on a en effet
\diagram{\begin{tikzcd}[ampersand replacement=\&,column sep=1em,row sep=1.1em]
\&1\arrow[d]\\
1\arrow[r]\&G'\arrow[r]\arrow[d]\&G\arrow[r]\&{\unsure{\mathfrak S|\mathfrak G}_{\tau}}\arrow[d,equal]\arrow[r]\&1\&{\big)}\\
\&{\displaystyle\mathop{\deleted{\mathfrak S}\mkern-14mu\prod}\limits_{i\in\tau}\mathfrak S_{S_i}}\arrow[d,"{\prod_{i\in\tau}\mathrm{sg}_i}"]\&\&{\mathbb Z/2\mathbb Z}\\
\&{\mathbb Z/2\mathbb Z}\arrow[d]\\
\&1
\end{tikzcd}}

% L16–17
\deleted{Le sous-groupe}

Une structure hexagonale subordonnée à la
% archive sheet number: 3

